\documentclass{article}
\usepackage{graphicx} % Required for inserting images
\usepackage{hyperref}
\usepackage{xcolor}

\title{DS-CV\&IC project of Group26}
\author{Evan F \& Kundai M}
\date{March 2026}

\begin{document}

\maketitle

Deliver at: \href{https://maat.eemcs.utwente.nl/team/239340}{maat.eemcs.utwente.nl}

\section{Motivation}
% \textcolor{gray}{Describe why your project with its questions and contributions is relevant and important. If you use literature, please provide references. 300 words max}

During disasters, relief organizations want to be able to help people in need as soon as possible. Information is a critical factor in aiding people in the right place at the right place. One method that can be used to assist the relief organizations is to use visual data from satellite images or posts from social media made by people on location. To actually get information from these images these need to be labelled. When this is done by people this takes a lot of time and labour. In the last years a lot of progression has been made in the field of computer vision to perform this task automatically. Now there are a lot of different models that are able to recognize images. We have chosen 5 models that are fit of doing image recognizion and in this paper we describe our method of testing different of these models on there performance on training and recognizing different types of disasters. We have done this by developing an automated 12-class disaster image classifier on a dataset of 7,303 images, where class imbalance and limited labeled data make robust modeling challenging. To keep training efficient while maintaining performance, we use transfer learning in a pipeline, then evaluated the reliability through 5-fold cross-validation and a separate testing set.
The result of a high accuracy could aid the relief organizations to map a disaster area using images, their disaster type labels and corresponding location contained in the meta data of these images.

% \begin{itemize}
%     \item Explain why automated disaster image classification is useful for crisis response and resource allocation.
%     \item State that manual labeling is slow and error-prone when many images arrive quickly.
%     \item Mention that transfer learning reduces training cost when labeled data is limited.
%     \item Describe your contribution: a ResNet34-based pipeline with cross-validation and hold-out testing.
%     \item Briefly note expected impact: faster and more consistent situational awareness.
% \end{itemize}

\section{(Business/Research) questions}
\textcolor{gray}{Specify the questions that your project has answered. Use the SMART criteria. 300 words max}

Our primary question that we want to be able to answer in this paper is how can we identify the most reliable transfer-learning image classifier for 12 disaster categories? Done by comparing five pre-trained models on the 7,303-image Incidents1M dataset using stratified 5-fold cross-validation and a final testing set. As a primary success metric we will use a weighted F1 score and macro F1, balanced accuracy, and recall as supporting metrics.

To answer this primary question we will answer the following subquestion:
Which of the five pre-trained models (ResNet50, VGG16, DenseNet121, MobileNetV3-Large, ViT-B/16) achieves the best cross-validated weighted F1?
How stable is each model across folds (mean and standard deviation), and which model is most robust rather than just best on one split?
To what extent do imbalance-aware techniques (class-weighted loss and weighted sampling) improve minority-class behavior (macro F1, balanced accuracy, recall)?
Does the selected model’s performance on the unseen hold-out test align with cross-validation results?
Which disaster classes are most frequently confused, and what does this imply for operational use in crisis-response decision support?

% \begin{itemize}
%     \item Main question: Can transfer learning classify 12 disaster-related image classes with reliable performance?
%     \item Sub-question: Which metric is most informative for this imbalanced multi-class problem (weighted F1 vs accuracy)?
%     \item Sub-question: How stable is model performance across folds?
%     \item Sub-question: Does hold-out test performance align with cross-validation estimates?
%     \item Write each question in SMART form: specific task, measurable metric, feasible scope, course-relevant, and deadline-bound.
% \end{itemize}

% Specific: Targeting a particular area for improvement
% Measurable: Quantifying, or at least suggesting, an indicator of progress
% Assignable: Defining responsibility clearly
% Realistic: Outlining attainable results with available resources
% Time-related: Including a timeline for expected results

% Specific: Which pretrained vision model best classifies 12 disaster/accident image classes for this dataset?
% Specific: How much does transfer learning with frozen backbones + new classifier head improve class-level recognition versus relying on a single metric?
% Measurable: Compare models using weighted F1, macro F1, balanced accuracy, accuracy, and recall across folds.
% Achievable: Restrict candidates to five widely used pretrained architectures and fixed training budget (epochs, batch size, optimizer, scheduler).
% Relevant: Accurate disaster-type recognition supports faster triage, monitoring, and downstream decision support.
% Time-bound: Complete model comparison with 5-fold CV, then retrain selected model and evaluate once on held-out test data.


\begin{itemize}
    \item Main question: Can transfer learning classify 12 disaster-related image classes with reliable performance?
    \item Sub-question: How stable is model performance across folds?
    \item Sub-question: Does hold-out test performance align with cross-validation estimates?
    \item Sub-question: 
    \item Write each question in SMART form: specific task, measurable metric, feasible scope, course-relevant, and deadline-bound.
\end{itemize}

\section{Source data}
\textcolor{gray}{What data quality issues did you consider that could potentially hamper your project (if any), what actions did you do to detect them (if any), what data quality issues did you find (if any), and what did you do about them? 500 words max}

\begin{itemize}
    \item Describe dataset source, number of classes, and examples of class names.
    \item Report class distribution and show that classes are imbalanced.
    \item Mention potential quality issues: mislabeled images, duplicate images, low-resolution/noisy images, and class overlap.
    \item Explain detection actions: visual inspection, class count plots, and stratified splitting checks.
    \item Explain mitigation: weighted sampler, class-weighted loss, and stratified train/val/test splits.
    \item State remaining risks that were not fully solvable (for example hidden label noise).
\end{itemize}

\section{Method}
\textcolor{gray}{Describe the main steps in your project with associated argumentation why they are relevant and important. Include methodically relevant details such as evaluation metrics and other methodically important parameters. Validation Method: One important aspect of modelling is validation. Especially for topics, DM, CV\&IC, GIS, and IENLP explicitly mention which validation method is used.
CV\&IC: Explicitly describe your approach to representing images and provide reasoning for your chosen representation. Include in the general question on Method explicitly which ML methods you have compared and how you compared them including which metrics you used for the comparison. 500 words max}

\begin{itemize}
    \item Summarize workflow: loading data, preprocessing, splitting, model setup, training, validation, and testing.
    \item Explain image representation: RGB tensors resized and cropped to 224x224, normalized with ImageNet statistics.
    \item Explain augmentation choices for training: random horizontal flip and small rotations.
    \item Describe model choice: pretrained ResNet34 with replaced final fully connected layer for 12 classes.
    \item Describe optimization setup: SGD, learning-rate scheduler, weighted cross-entropy, batch size, and epoch counts.
    \item Describe validation method explicitly: Stratified 5-fold cross-validation on training data plus final hold-out test.
    \item Describe comparison approach: compare fold-wise weighted F1, accuracy, precision, recall, and confusion matrices.
    \item Justify weighted F1 as key metric due to class imbalance.
\end{itemize}

\section{Results}
\textcolor{gray}{Provide tables, graphs, and visualizations for your results with associated descriptions how to interpret them and which detailed conclusions you draw from them (text in the tables, graphs, and visualizations or in their captions do not count for the word limit) 500 words max}

\begin{itemize}
    \item Report mean and variability of cross-validation weighted F1 and accuracy.
    \item Add learning-curve plots (train/validation loss and accuracy) and explain overfitting signs or stability.
    \item Add fold-wise F1 plot and interpret consistency across folds.
    \item Include confusion matrices (fold-level and final test) and discuss frequent confusions between classes.
    \item Report hold-out test metrics: accuracy, precision, recall, weighted F1, and class-wise report.
    \item Conclude whether test results are close to CV estimates and what this means for robustness.
\end{itemize}

\section{Reliability of results}
\textcolor{gray}{Describe limitations and explain how reliable your results and conclusions can be trusted to be. Describe any reliability-related actions (if any). 300 words max}

\begin{itemize}
    \item Discuss reliability strengths: stratified splits, k-fold validation, and separate hold-out test set.
    \item Discuss limitations: limited dataset size, possible label noise, and domain shift risk.
    \item Mention randomness sources: initialization, data loader sampling, augmentation effects.
    \item State reliability actions: fixed random seed, repeated evaluation across folds, and metric triangulation.
    \item Clarify that conclusions are reliable within similar data conditions, but external deployment needs extra validation.
\end{itemize}

\section{Technical depth}
\textcolor{gray}{What methods/techniques did you apply in the project that go beyond the primary topic and where did you learn about them (other topic, literature, other course, etc.) 300 words max}

\begin{itemize}
    \item Highlight transfer learning with pretrained CNNs as core advanced technique.
    \item Highlight imbalance handling using weighted sampler and class-weighted loss.
    \item Highlight multi-level evaluation design: fold metrics, learning curves, confusion matrices, and hold-out verification.
    \item Mention practical deep learning choices: scheduler tuning, augmentation policy, and checkpointing best model.
    \item State where techniques were learned (course material, PyTorch tutorial, and relevant literature).
\end{itemize}

\section{Conclusions \& recommendations}
\textcolor{gray}{Describe the target stakeholder, the main conclusions you drew from your results, the recommendations for this target stakeholder, and future research. 300 words max}

\begin{itemize}
    \item Identify stakeholders: emergency management teams, monitoring centers, or public safety analysts.
    \item Conclude whether the model is suitable for decision support and where it performs best/worst.
    \item Recommend deployment as an assistive tool, not a fully autonomous decision-maker.
    \item Recommend collecting more balanced data for weak classes and adding label-quality checks.
    \item Suggest future work: stronger backbones, ensembling, calibration, and external dataset validation.
\end{itemize}

\section{Reflection}
\textcolor{gray}{How did the project’s challenges get addressed, how relevant were the course’s skills and techniques for this real-world project, and what additional knowledge and skills may have been beneficial? For transparency, did you use ChatGPT and if so, for what purpose(s)? 300 words max}

\begin{itemize}
    \item Describe key project challenges: imbalance, multi-class confusion, and balancing accuracy with generalization.
    \item Explain how course skills helped: preprocessing, model validation, metric interpretation, and result communication.
    \item Mention additional skills that would help: error analysis pipelines, experiment tracking, and model explainability.
    \item If AI tools were used, state exactly how (for editing text, debugging, or structure support) and what you verified yourself.
    \item Reflect on what you would do differently in a second iteration.
\end{itemize}

\section{ML issues}
\textcolor{gray}{CV\&IC: Describe which ML issues you addressed (such as class imbalance, etc.) and how you addressed them. 200 words max}

\begin{itemize}
    \item Describe class imbalance and how you addressed it with weighted sampling and weighted loss.
    \item Mention overfitting risk and your controls: augmentation, validation monitoring, and best-epoch model selection.
    \item Mention metric bias risk from imbalance and why weighted F1 is emphasized over plain accuracy.
    \item Mention potential data leakage risk and how strict split strategy reduced it.
\end{itemize}

\section{Generalisation capabilities of the method}
\textcolor{gray}{CV\&IC: Describe how you evaluate the performance of your model and its generalisation ability. 200 words max}

\begin{itemize}
    \item Explain that generalization was assessed using stratified 5-fold cross-validation and a final unseen test set.
    \item Report both average performance and fold variability, not only a single best score.
    \item Compare CV mean F1 with hold-out test F1 to assess stability.
    \item Use confusion matrices and class-wise metrics to show where generalization is strong or weak.
\end{itemize}

\section{References}
\textcolor{gray}{(Optional) Add a list of references used in your project. 1000 words max}

\bibliography{ref.bib}


https://github.com/onnx/models?tab=readme-ov-file#image_classification
https://ieeexplore.ieee.org/document/10134069

\section{Appendix}
\textcolor{gray}{(Optional) Add appendix content to your report. 1000 words max}

\end{document}
